% The outline of the course (Ch1_1, Ch2_1, Ch3_1, Ch4_1 slide 2).  \ELouton = 1 | 2 | 3 | 4: the chapter the
% slide highlights.  A text hierarchy:
% the root stands at the leading edge, so it is mirrored between the editions.
\providecommand{\ELouton}{1}
\begin{tikzpicture}[x=\ELrsgn cm,y=1cm]
  \def\ELon#1#2{\ifnum\ELouton=#1 #2\fi}
  \node[ELtreeroot,text width=2.9cm] (root) at (6.9,0) {\ELL{Engineering Electromagnetics}{الکترومغناطیس مهندسی}};
  \node[ELtreebox,text width=3.9cm,\ELon{1}{ELtreeon}] (m1) at (2.6,4.05) {\ELL{Familiarity with and mastery of the required mathematical tools}{آشنایی و تسلط بر ابزار ریاضی مورد نیاز}};
  \node[ELtreeboxR,text width=3.9cm,\ELon{2}{ELtreeonR}\ELon{3}{ELtreeonR}] (m2) at (2.6,1.6) {\ELL{The effect of static electric charges}{بررسی اثر بارهای الکتریکی ساکن}};
  \node[ELtreebox,text width=3.9cm,\ELon{4}{ELtreeon}] (m3) at (2.6,-1.4) {\ELL{The effect of electric charges moving with constant velocity}{بررسی اثر بارهای الکتریکی متحرک با سرعت ثابت}};
  \node[ELtreeboxR,text width=3.9cm] (m4) at (2.6,-4.05) {\ELL{The effect of electric charges moving with varying velocity}{بررسی اثر بارهای الکتریکی متحرک با سرعت متغیر}};
  \node[ELtreebox,text width=4.3cm,\ELon{1}{ELtreeon}] (c1) at (-2.2,4.05) {\ELL{Chapter 1: Vector Analysis}{فصل اول: آنالیز برداری}};
  \node[ELtreeboxR,text width=4.3cm,\ELon{2}{ELtreeonR}] (c2) at (-2.2,2.3) {\ELL{Chapter 2: Static Electric Fields}{فصل دوم: میدان‌های الکتریکی ساکن}};
  \node[ELtreeboxR,text width=4.3cm,\ELon{3}{ELtreeonR}] (c3) at (-2.2,0.9) {\ELL{Chapter 3: Solution of Electrostatic Problems}{فصل سوم: حل مسائل الکتریسیته ساکن}};
  \node[ELtreebox,text width=4.3cm,\ELon{4}{ELtreeon}] (c4) at (-2.2,-0.7) {\ELL{Chapter 4: Steady Electric Currents}{فصل چهارم: جریان‌های الکتریکی دائم}};
  \node[ELtreebox,text width=4.3cm] (c5) at (-2.2,-2.1) {\ELL{Chapter 5: Static Magnetic Fields}{فصل پنجم: میدان‌های مغناطیسی ساکن}};
  \node[ELtreeboxR,text width=4.3cm] (c6) at (-2.2,-4.05) {\ELL{Chapter 6: Time-Varying Electromagnetic Fields}{فصل ششم: میدان‌های الکترومغناطیسی متغیر با زمان}};
  \foreach \m in {m1,m2,m3,m4}{\draw[ELtreelineG] (root.\ELend)--(\m.\ELstart);}
  \draw[ELtreeline] (m1.\ELend)--(c1.\ELstart);
  \draw[ELtreelineR] (m2.\ELend)--(c2.\ELstart) (m2.\ELend)--(c3.\ELstart);
  \draw[ELtreeline] (m3.\ELend)--(c4.\ELstart) (m3.\ELend)--(c5.\ELstart);
  \draw[ELtreelineR] (m4.\ELend)--(c6.\ELstart);
\end{tikzpicture}
